
Vladimir Gusev~\cite{Gusev2012} introduced the constrained synchronization problem
in connection with a reduction procedure on automata
that preserves the reset length up to plus/minus one w.r.t.~a given constraint language. 
We discuss this in \cref{constrsync:sec:motivation:reduction}.
The problem of constrained synchronization also appears implicitly 
in a reduction used by Martyugin~\cite{Martyugin14}. %, for which we give details below
Here, we give details on the procedure from~\cite{Gusev2012} and on
the reduction used in~\cite{Martyugin14} in \cref{constrsync:sec:motivation:implicit}.
Then we outline a relation to the \emph{Skolem problem}~\cite{BLONDEL2002,DBLP:conf/rp/OuaknineW12} and the \emph{matrix orbit problem}~\cite{KannanLipton1986} in \cref{constrsync:sec:motivation:skolem}.
Finally, we sketch possible applications that naturally extend
areas where synchronizing automata have been used previously
in \cref{constrsync:sec:motivation:operating,constrsync:sec:motivation:part-orienters,constrsync:sec:motivation:supervisory}.

\section{Use in a Reduction Procedure}
\label{constrsync:sec:motivation:reduction}

\IN{$r$-decomposition}
\IN{rank factorization}
Let $\semiautomat{A}{Q}{\Sigma}{\delta}$ be a semiautomaton with $n$ states
over a binary alphabet $\Sigma = \{a,b\}$. A set of matrices $X_1, X_2, Y_1, Y_2$
over $\{0,1\}$ %Boolean semiring
is an \emph{$r$-decomposition}~\cite{Gusev2012} of $\mathcal A$
if
\begin{enumerate}
\item $X_1$ and $Y_1$ are $n \times r$ matrices,
\item $X_2$ and $Y_2$ are $r \times n$ matrices,
\item $[a] = X_1 X_2$, $[b] = Y_1 Y_2$,
\item ev ery row of $X_1, X_2, Y_1, Y_2$ has only one occurrence of~$1$.
\end{enumerate}
The automaton $\mathcal A$ admits an $r$-decomposition if and only if $r \ge \max\{\rank{a},\rank{b}\}$~\cite{Gusev2012}.
By the first three conditions $X_1, Y_1$ and $X_2,Y_2$ constitute
\emph{rank factorizations}
of $[a]$ and $[b]$, respectively, as discussed in~\cite{piziak99}.
Since the matrices $X_1, X_2, Y_1, Y_1$ have one~$1$ in each row,
the same holds for their products.
Hence, we obtain a complete and deterministic $r$-state automaton
$\semiautomat{A_{X_1, X_2, Y_1, Y_2}}{Q_{X_1, X_2, Y_1, Y_2}}{\Gamma}{\delta_{X_1,X_2, Y_1, Y_2}}$
by setting $Q_{X_1, X_2, Y_1, Y_2} = \{1,\ldots,r\}$ and $\Gamma = \{ a_{X_2Y_1}, a_{X_2X_1}, a_{Y_2X_1}, a_{Y_2Y_1} \}$
where the action of the letters is given
by $[a_{X_2Y_1}] = X_2 Y_1$, $[a_{X_2X_1}] = X_2X_1$, $[a_{Y_2X_1}] = Y_2 X_1$
and $[a_{Y_2 Y_1}] = Y_2 Y_1$. % (w.r.t.~some fixed ordering of the states $Q_{X_1, X_2, Y_1, Y_2}$, see \cref{prelim:subsec:transmtx}).
The following was shown by Gusev~\cite[Prop.~5]{Gusev2012}.

\begin{proposition}
		 Let $\semiautomat{A}{Q}{\{a,b\}}{\delta}$ be an automaton with $r$-decomposition $X_1, X_2, Y_1, Y_2$
		 and $\semiautomat{A_{X_1, X_2, Y_1, Y_2}}{Q_{X_1, X_2, Y_1, Y_2}}{\Gamma}{\delta_{X_1,X_2, Y_1, Y_2}}$.
		 Then $\mathcal A$ is synchronizing if and only if 
		 \[ Syn(\mathcal A_{X_1, X_2, Y_1, Y_2}) \cap R \ne \emptyset
		 \] 
		 with
		 \begin{align*}
		 	R = \Gamma^* \setminus \big(
		 	& \Gamma^* ( a_{X_2 X_1} ( a_{Y_2 X_1} + a_{Y_2 Y_1} ) ) \Gamma^* + 
		 	\Gamma^* ( a_{X_2 Y_1} ( a_{X_2 Y_1} + a_{X_2 X_1} ) ) \Gamma^* + \\ 
		 	& \Gamma^* ( a_{Y_2 Y_1} ( a_{X_2 Y_1} + a_{X_2 X_1} ) ) \Gamma^* + 
		 	\Gamma^* ( a_{Y_2 X_1} ( a_{Y_2 X_1} + a_{Y_2 Y_1} ) ) \Gamma^* \big),
		 \end{align*}
	     where, if this is the case, $m - 1 \le \rt{\mathcal A} \le m + 1$
	     for  $m = \min\{ |u| : u \in \Syn(\mathcal A_{X_1, X_2, Y_1, Y_2}) \cap R  \}$.
\end{proposition}

So, the reset length of an automaton $\mathcal A$ without permutational letters can be bounded by the length of a shortest synchronizing word 
of a smaller automaton $\mathcal A_{X_1, X_2, Y_1, Y_2}$ contained in a constraint language independent of $\mathcal A$.

Gusev~\cite{Gusev2012} used $r$-decompositions to construct $n$-state automata 
with reset length at least $r^2 - r - 1$. This was the sole purpose of this work
and while the above statement used constrained synchronization, it was not introduced
as a decision problem and questions of computational complexity were not 
investigated.

\section{Implicit in a Reduction}
\label{constrsync:sec:motivation:implicit}

Martyugin~\cite{Martyugin14} 
proved that it is \PSPACE-complete to decide whether a given partial automaton
has a carefully synchronizing word (for the definition of these notions, see \cref{constr_sync:subsec:pdfa}).
His reduction constructs a partial automaton over $\Sigma$
where a shortest carefully synchronizing 
word has the form $x\Gamma^*y$ for $x,y \in \Sigma$ and $\Gamma = \Sigma \setminus \{x,y\}$~\cite[Lem.~1]{Martyugin14}.
By adding self-loops to obtain 
a complete automaton in the construction,
one can easily see that it has a synchronizing word in $x\Gamma^*y$
if and only if the original partial automaton
has a carefully synchronizing word.
In this sense, \constrsync{x\Gamma^*y} is implicit
in the construction. 

\section{The Matrix Orbit Problem and the Skolem Problem}
\label{constrsync:sec:motivation:skolem}

\IN{Skolem problem} 
\IN{Pisot's problem} 
\IN{matrix orbit problem}
The \emph{(matrix) orbit problem} is the problem
which asks for an $n \times n$ matrix $A$ over $\mathbb Q$
and
$x,y \in \mathbb Q^n$ %a matrix $A \in \mathbb Q^{n \times n}$
whether $xA^i = y$
for some $i \ge 0$. This problem was shown to be solvable
in polynomial time by Kannan \& Lipton~\cite{KannanLipton1986}.
In~\cite[Cor.~1.3]{BLONDEL2002}
the problem which asks for an $n \times n$ matrix $A$ over~$\mathbb Z$
and $1 \times n$ and $n \times 1$ vectors $x$ and $y$ over $\mathbb Z$, respectively, 
whether $xA^iy = 0$ for some $i \ge 0$ was shown to be \NP-hard.
As noted in~\cite{BLONDEL2002},
this implies \NP-hardness of the \emph{Skolem problem} (also called \emph{Pisot's problem})
which ask whether an integer linear recurrent sequence has zero as a value,
see \cite{BLONDEL2002}, and also~\cite{DBLP:conf/rp/OuaknineW12} for more information on the Skolem problem.
Decidability of the Skolem problem is unknown~\cite{DBLP:conf/rp/OuaknineW12}.

Next, we point out a connection
to the two problems above.\footnote{These connections were originally pointed out by Vladimir V. Gusev (personal communication).}
Let $\semiautomat{A}{Q}{\{a,b\}}{\delta}$ with $Q = \{q_1,\ldots,q_n\}$
and $q \in Q$.
Set
\[
V_q = \{ (y_1, \ldots, y_n) \mid y_j > 0 \text{ implies } q_j \in \delta^{-1}(q, a) \text{ for } j \in \{1,\ldots,n\} \},
\]
and define numbers $x_1, \ldots, x_n$ and $z_1, \ldots, z_n$
by 
\[
x_i = \begin{cases}
	1 & \text{if } q_i \in \delta(Q, a), \\
	0 & \text{otherwise,}
\end{cases}
\quad\text{ and }\quad
z_i = \begin{cases}
	1 & \text{if } \delta(q_i, a) \ne q, \\
	0 & \text{otherwise,}
\end{cases}
\]
for $i \in \{1,\ldots,n\}$. Then for $i \ge 0$ we have
\[
	\delta(Q, ab^ia) = \{q\}
	 \Leftrightarrow x \cdot [b]_{\mathcal A}^i \in V_q 
	 \Leftrightarrow  x \cdot [b]_{\mathcal A}^i \cdot \begin{pmatrix} z_1 \\ \vdots \\ z_n \end{pmatrix} = 0
\]
We will show later in \crefitem{constr_sync:thm:introductory_examples}{constrsync:ex:item:abstara}
that $\constrsync{ab^*a}$ is \NP-complete. The above construction
then also yields \NP-hardness of the problem introduced above
that asks whether $xA^iy = 0$ for a matrix~$A$ and vectors~$x,y$ over~$\mathbb Q$.
The construction also yields another proof
that $\constrsync{ab^*a}$ is in \NP\  since checking $x \cdot [b]_{\mathcal A}^i \in V_q$ for some $q \in Q$
is in \NP\  by guessing $q \in Q$ and then checking $y \in V_q$ and $x[b]^i_{\mathcal A} = y$,
which can be done in polynomial time since the matrix orbit problem 
can be solved in polynomial time.
 
\begin{remark}
	In \cref{prelim:compl:thm:reach-comm-fixed} 
	we proved that \REACHF
	is solvable in polynomial time for commutative automata over a fixed alphabet.
	The result of Kannan \& Lipton~\cite{KannanLipton1986}
	stating polynomial-time solvability 
	of existence of some $i \ge 0$
	such that $xA^i = y$ for a matrix~$A$
	and vectors $x,y$ over~$\mathbb Q$
	also implies polynomial-time solvability
	of \REACHF for unary automata
	similarly as above by considering the transition matrix and encoding state
	sets $S$ into vectors $(x_1, \ldots, x_n)$ 
	by setting $x_i = 1$ if $q_i \in S$ and $x_i = 0$ otherwise.
	Note that the related problem \ST 
	is \NP-complete for unary automata by \cref{compl:thm:comm_ST}.	
\end{remark}

\section{Operating Restrictions}%ursprungsbeispiel, aber auch homogene experimente sternberg paper
\label{constrsync:sec:motivation:operating}

Fernau et al.~\cite{FernauGHHVW19} gave the following motivating example. Suppose that a system, modeled by an automaton, 
has a ``normal'' and a ``directing`` operating mode, and only in the latter mode
synchronizing instructions are allowed. The switching between both modes 
can be modeled, in its simplest form, 
by the constraint $ab^*a$, where $a$ switches into the ``directing'' mode.
As written in the previous paragraph, \constrsync{ab^*a} is \NP-complete, see \crefitem{constr_sync:thm:introductory_examples}{constrsync:ex:item:abstara}.
Even more generally, it might be possible that a system, for example, a remotely controlled
rover on a distant planet or a satellite in orbit, is not allowed to execute all commands in every possible
order. For example, some commands are only allowed in a certain order or after
other commands have been executed. All of this imposes constraints
on the possible reset sequences.

% bestimmte operationen nur wenn bestimmte andere, oder erst vertikal
\section{Part Orienters}
\label{constrsync:sec:motivation:part-orienters}

Generalizing the part orienters example
from \cref{intro:sec:examples:partorienters}, suppose some parts arrive at a manufacturing site and they
need to be sorted and oriented before assembly. Practical considerations
favor methods which require little or no sensing, employ simple devices,
and are as robust as possible. 
As in \cref{intro:sec:examples:partorienters},
the parts are put on a conveyor belt
so that the stream of the parts encounter a series of passive obstacles
placed along the belt before reaching the assembly point.
Now, furthermore, assume, due to the availability of components, requirements how to lay things out
or physical restrictions,
the passive components can %could 
not be placed at random along
the belt, but have to obey some restrictions, or restrictions
in what order they are allowed to happen. 
These constraints on the
permissible sequences of applicable operations 
can be modeled by the constrained synchronization problem.


\IN{discrete event system}
\section{Supervisory Control} 
\label{constrsync:sec:motivation:supervisory}

Ramadge and Wonham~\cite{RamadgeWonham87} modelled
 a \emph{discrete event system} (DES) by a partial automaton 
 generating an event sequence (see also \cite{DBLP:books/daglib/0034521,wonham2019}). 
 %The authors distinguished between controllable and uncontrollable
 %events.
 Then, a supervisor,  depending on the event sequence \todo[inline]{supervisor or controller}
 that occurred so far,
 is able to restrict the set of events
 that are possible in the next step, where, however,
 he can only limit the use of so-called controllable events.
 So, if we want to (globally) reset a finite state DES~\cite{Alves2020} under supervisory control, 
 this can be seen as an instance of a constrained synchronization problem
 in case the permissible event sequences form
 a regular language.
 
 

% % noch etwas umschreiben, sonst identsich mit Mischa-Text
% \item[Biocomputing.] In~\cite{Benenson2003,Benenson2001} DNA molecules have been used as both
% hardware and software for finite automata of nanoscaling size, see also~\cite{Vol2008}. 
% For instance, Benenson et al~\cite{Benenson2003} produced ``a `soup of automata', that is, a solution
% containing $3 \times 10^{12}$ identical automata per $\mu1$. All these molecular
% automata can work in parallel on different inputs, thus ending up in different
% and unpredictable states. In contrast to an electronic computer, one cannot
% reset such a system by just pressing a button; instead, in order to synchronously bring
% each automaton to its start state, one should spice the soup with (sufficiently
% many copies of) a DNA molecule whose nucleotide sequences encodes a reset word''~\cite{Vol2008}.
% Now, it might be possible that certain sequences, or subsequences,
% are not possible as they might have unwanted biological side-effects, 
% or might destroy the molecules at all.
